\begin{problem}{Jeweler's Fair Cuts}{standard input}{standard output}{2 seconds}{256 megabytes}

Master jeweler Jade has acquired a necklace of $n$ gemstones strung in a row; gem $i$ is appraised at $a_i$ ducats --- a few \emph{debt stones} among them are actually IOUs, so appraisals may be negative.

For the guild fair, Jade must split the row into one or more contiguous strands so that every gem belongs to exactly one strand. A split is called \emph{balanced} if all its strands carry the same total appraisal.

Count the number of balanced splits. Two splits are considered different if their sets of cut positions are different (a cut can be made between positions $i$ and $i+1$ for any $1 \le i < n$).

Note that selling the necklace whole --- a single strand, no cuts --- is always balanced.

Since the number of splits can be huge, print it modulo $10^9 + 7$.


\InputFile
The first line contains a single integer $t$ ($1 \le t \le 10^5$) --- the number of test cases.

The first line of each test case contains a single integer $n$ ($1 \le n \le 10^5$) --- the size of the array.

The second line of each test case contains $n$ integers $a_1, a_2, \ldots, a_n$ ($-10^9 \le a_i \le 10^9$) --- the appraisals of the gems, in necklace order.

It is guaranteed that the sum of $n$ over all test cases does not exceed $10^5$.


\OutputFile
For each test case, print a single integer --- the number of balanced splits of the necklace, modulo $10^9 + 7$.


\Example

\begin{example}
\exmpfile{example.01}{example.01.a}%
\end{example}

\Note
The example contains six test cases.

\begin{itemize}
  \item $a = [1, 2, 1, 2]$, total sum $6$. Balanced divisions: the whole array (one segment of sum $6$), and $[1, 2]\;[1, 2]$ (two segments of sum $3$). Three segments of sum $2$ would require a prefix of sum $2$, which does not exist. Answer: $2$.
  \item $a = [1, -1, 1, -1]$, total sum $0$, so every segment must sum to $0$. The only position $i < n$ with prefix sum $0$ is $i = 2$, and we may either cut there or not. Answer: $2$.
  \item $a = [1, 1]$: $[1, 1]$ or $[1]\;[1]$. Answer: $2$.
  \item $a = [5]$: only the single segment. Answer: $1$.
  \item $a = [-1, -1]$: same as the third case with a negative total. Answer: $2$.
  \item $a = [1, 2, 3, 0, 2, 4]$, total sum $12$. With $k = 2$ segments (sum $6$ each) the cut can be made after index $3$ or after index $4$, since both prefix sums equal $6$ --- that gives $2$ divisions. No $k \in \{3, 4, 6\}$ works (e.g. for $k = 3$ a prefix of sum $4$ is required but absent). Together with the single-segment division, the answer is $3$.
\end{itemize}


\end{problem}

